\section{Conclusion}
\label{sec:conclusion}

PSBD transfers to a Vision Transformer once its perturbation is placed where a ViT carries a backdoor. Our reading of the original placement, dropout on the residual stream, detects the attacks with a firm shortcut and misses the rest, and it ranks \DetectorsAurocRankPublished{} of the \DetectorsAurocDefensesRanked{} defenses we compare. Moving the same dropout to the attention input repairs part of that. Masking whole tokens there, which we call PSBD-TM, reaches a mean AUROC of \HeadlineAurocAdaptive{} on \PanelCellsCached{} ViT-B/16 models and \SwinRecommendedAurocAdaptive{} on \SwinCells{} Swin-S models. It ranks \DetectorsAurocRankOurs{} of \DetectorsAurocDefensesRanked{} and gains most at \PanelRateLow{} poisoning, where a detector is needed most.

The placement is not a tuning accident. A backdoor in a ViT is 1 direction in the residual stream, written in the last third of the network and routed through attention from the trigger's patch tokens to the class token. A perturbation that removes tokens before attention removes the shortcut, while a perturbation on the residual stream perturbs the decision itself. The same account explains why noise and masking swap places across a LayerNorm, predicts which attacks show the prediction shift toward the target class that the original paper observed, and reproduces the direction removal of Karayal\c{c}{\i}n et al.\ once the direction is removed from every residual write.

For a practitioner the recommendation is short. On a ViT or a Swin, run PSBD with token masking at the attention input and choose the rate by the clean shift ratio as the original does. Against a knowledgeable attacker, combine several placements, since an attacker who evades 1 probe does not evade the others. The union of probes costs 1 sweep per probe and restores a mean AUROC of \AdaptiveVitUnion{} on the models trained to evade a single probe.

We also state what the result is not. Against the strongest competitor, the calibrated IBD-PSC port, we are \DetectorsAurocMargin{} ahead on AUROC and \DetectorsTprTwoZeroMargin{} behind at a 20\% false-positive budget. That is a tie. The contribution is that the placement decides whether prediction-shift detection works on this architecture at all, not that the method now leads the field.

2 questions remain open. An attacker who trains against the union of probes rather than a single one. And the all-to-all attack, whose backdoor has no single direction and which, on the present evidence, needs a different statistic.
